% Source: 36_Chapter_26_Thrust_Into_the_Game_John_Calvin.txt
\bookchapter[Thrust into the Game]{CHAPTER 26}{THRUST INTO THE GAME}{John Calvin}

\leadin{T}{HE WAR RAGING BETWEEN} Spain and France had closed the road to Strasbourg. So the young
scholar from France, hoping to continue his studies, was forced to go through Geneva. He planned to
stay only one night. He knew that the town wasn’t the best place for solitude.

Geneva was in disarray. It was notoriously pleasure-loving, but its recent rejection of the Duke of
Savoy and the pope in Rome had left public affairs in shambles, torn by dangerous factions.

An inflammatory Reformer named William Farel had been preaching in Geneva for four years, and
Catholic masses had ceased. But Geneva’s Protestantism rested chiefly on political hostility to the
bishop, not doctrinal convictions. No one had stepped forward to shape the city’s religious
institutions along biblical lines.

Farel knew the city needed a manager. So during the “chance” visit of the young Frenchman John
Calvin, Farel made a point to call upon the visitor. He found in Calvin the answer to Geneva’s need,
and he urged the young scholar to go no farther but to stay in the city and help establish the work
there.

Calvin protested that he had some special studies he wanted to pursue. But Farel responded, “You are
only following your own wishes! If you do not help us in this work of the Lord, the Lord will punish
you for seeking your own interest rather than His.”

Calvin was terror-stricken. The last thing he wanted was to offend Almighty God! So he consented to
stay and immediately took up the reforming cause in Geneva.

Years later, looking back over his career, Calvin observed, “Being by nature a bit antisocial and
shy, I always loved retirement and peace~.~.~.~. But God has so whirled me around by various events
that He has never let me rest anywhere, but in spite of my natural inclination, has thrust me into
the limelight and made me ‘get into the game,’ as they say.”

Calvin’s leadership in “the game” shaped a third reformation tradition. Today we call it Reformed or
Calvinistic Christianity. It includes all Presbyterians, Dutch and German Reformed Churches, and
many Baptists and Congregationalists.

\section{THE MEANING OF REFORMED CHRISTIANITY}

What are the unique features of Reformed Christianity? We can trace the dominant characteristics to
the life and teachings of that young scholar who heard the call of God in William Farel’s warning.

God had a game. Calvin was convinced of it. He called it God’s sovereign will. Just as Luther’s
central doctrine was justification by faith, so Calvin’s was the sovereignty of God. Both Reformers
had an overwhelming sense of the majesty of God, but Luther’s served to point up the miracle of
forgiveness, while Calvin’s gave the assurance of the impregnability of God’s purpose.

Calvin (1509–1564) shared with Luther the four central Protestant beliefs, but he was born a
generation after Luther in a different land and was a far different sort of person.

Luther was a peasant, a monk, and a university professor; Calvin, a scholar and lawyer called to a
turbulent public ministry in a flourishing business community. Naturally they were impressed by
different needs and emphasized different Christian solutions. Their foundations were the same, but
the structures of doctrine and practice they erected at Wittenberg and Geneva were different in many
important respects.

Calvin’s organizing and executive abilities enabled him to build on the work of Zwingli. The reform
movement started at Zurich spread rapidly in German-speaking Switzerland. Bern and its surrounding
territories (the canton) joined the Protestant ranks in 1528. Basel followed in 1529. Soon
Appenzell, St. Gall, and Schaflhausen joined the movement.

The Zwinglian type of reformation found support also outside of Switzerland in the important German
city of Strasbourg. There Martin Bucer (1491–1551), surpassed in influence throughout Germany only
by Luther and his associate Melanchthon, was inclined to sympathize with Zwingli rather than with
Luther.

Controversy between the Protestant and Catholic cantons in Switzerland led, however, on October 11,
1531, to a battle at Kappel between Zurich and its Catholic neighbors. In the course of the battle
Zwingli lost his life. Thus the leadership of the reformation in Zurich fell to Heinrich Bullinger.
By the 1540s, however, Geneva, in French-speaking Switzerland, had emerged as the international
center of Reformed Christianity under Calvin’s disciplined hand.

John Calvin came from a small town sixty miles northeast of Paris. His father was anxious for his
son to have the advantages of a good education. Calvin entered the University of Paris at fourteen
and mastered not only a brilliant writing style but a skill in logical argument. In later years men
might not like what Calvin said, but they could not misunderstand what he meant. He left the
university in 1528 with his Master of Arts degree.

After Paris, at his father’s insistence, John turned to the study of law in the universities of
Orleans and Bourges, but his father’s death in 1531 left Calvin free to pursue his own interests.
Thus he returned to Paris as a student of the classics, intent upon a scholar’s career.

His studies brought Calvin into touch with reforming ideas circulating in Paris, and shortly
thereafter an event in Calvin’s life turned him in a new direction. He called it an “unexpected
conversion.” We can’t be certain about the date, but it was clearly more than a kind of spiritual
enlightenment or the recognition of the supreme authority of the Scriptures. Calvin surrendered his
will to God. He gave up his career as a classical scholar and identified with the Protestant cause
in France.

In the autumn of 1533 Calvin was so closely linked with his friend Nicholas Cop that when Cop gave a
strongly Protestant address as rector of the university, some suspected Calvin wrote the speech. The
vigorous address threw the institution into an uproar, and Calvin was forced to flee from Paris. The
young Reformer found refuge in Basel where in March 1536 he published the first edition of his
highly influential Institutes of the Christian Religion.

The work was the clearest, most logical, and most readable exposition of Protestant doctrine that
the Reformation age produced, and it gave its youthful author European fame overnight. Calvin
labored on its elaboration nearly all his active life. Twenty years later it was a much larger work
but its interpretation of Christian truth remained essentially the same.

As a preface to the Institutes, Calvin addressed a remarkable letter to Francis I, King of France,
defending the Protestants in that land from the criticisms of their enemies and vindicating their
rights to a respectful hearing. No one had spoken so effectively in their behalf, and with this
letter Calvin assumed a position of leadership in the Protestant cause.

\section{GENEVA: HOME OF CALVINISM}

He could no longer remain in France. This is why after some months of travel he was headed for
Strasbourg that fateful night in July 1536 when William Farel enlisted him for the work of God in
Geneva.

The city councils offered Calvin a position, Professor of Sacred Scriptures, and he began his work
with vigor. He prepared a confession of faith to be accepted by everyone who wished to be a citizen,
he planned an educational program for all, and he insisted on excommunication, particularly
expulsion from the Lord’s Supper, for those whose lives did not conform to spiritual standards.

It was the most strenuous program of moral discipline within Protestantism, a bit more than the city
fathers had bargained for. Opposition arose, especially over who had the power to excommunicate, the
church or the magistrates. After a year of struggle Calvin and Farel were defeated on the issue, and
in April 1538 the city councils ordered the two Reformers to leave Geneva. Calvin appeared to be a
failure.

The next three years spent in Strasbourg were probably the happiest in Calvin’s life. As pastor of
the church of French refugees, he was free to carry out his disciplinary measures; he was a
successful teacher of theology; he was honored by the city and was made its representative to
important religious conferences in Germany. He married a widow with two children during these years
and she remained his helpful companion until her death in March 1549.

Meanwhile Calvin’s friends in Geneva regained power in the city government and urged him to return
and resume his reform efforts. In September 1541, with great reluctance, he once again took up the
Geneva burden, practically on his own terms. The church constitution accepted by the city put into
practice Calvin’s leading ideas.

Four offices governed the church: pastors, teachers, elders, and deacons. The twelve elders,
together with the ministers, formed the Consistory, which was responsible for the moral supervision
of the city. Offenses ranged from absences from public worship to drinking, adultery, gambling, and
dancing.

Naturally, opposition continued. Many times Calvin was on the brink of banishment. But he fought his
way through courageously, and the influx of exiles for their faith, chiefly from France, whom Calvin
attracted to Geneva constantly increased his following.

During one low point in Calvin’s influence, in 1553, the brilliant but erratic Spanish physician
Michael Servetus sought refuge in Geneva. Servetus was fleeing Catholic persecution for his heresy
of denying the doctrine of the Trinity. He arrived in Geneva just as Calvin’s enemies were
challenging his authority. While Calvin wanted a more merciful death than burning for the heretic,
he did support the silencing of the ill-balanced thinker. Servetus was burned at the stake, and many
in later generations remembered Calvin primarily as “the man who burned Servetus.”

Two years later Calvin’s position in Geneva was secure, and until his death he enjoyed significant
influence over religious and moral issues in the city. For Calvin, however, Geneva was never an end
in itself. He considered the city a refuge for persecuted Protestants, an example of a disciplined
Christian community, and a center for ministerial training. Enthusiastic students from all over
Europe came to Geneva to see what John Knox called “the most perfect school of Christ that ever was
on earth since the days of the Apostles.” They received Calvin’s theology firsthand and obediently
lived the city’s rigorous lifestyle.

\section{THE SOVEREIGNTY OF GOD}

The Calvinism the students imbibed flowed from Calvin’s central belief in the absolute sovereignty
of God. “God asserts his possession of omnipotence,” he wrote, “and claims our acknowledgement of
this attribute.” God is the “Governor of all things.” In his own wisdom, from the remotest eternity,
he decreed what he would do, and by his own power he executes what he has decreed.

This is more than a general guidance. The Bible teaches God’s particular direction in individual
lives. We read that not a sparrow falls to the ground unknown to the Father. We read, too, that he
has given babies to some mothers and withheld them from others. Such events represent no relentless
fatalism in nature, but the personal decrees of Almighty God, who moves men to walk in his ways.

If Luther’s ultimate text was “the just shall live by faith,” Calvin’s was, “Thy will be done, on
earth as it is in heaven.” Calvin saw the old doctrine of predestination taught by Paul, Augustine,
and Luther as a source of religious devotion. More than a problem of the mind, Calvin considered
divine election to eternal life the deepest source of confidence, humility, and moral power.

While Calvin did not profess to know absolutely who were God’s chosen (the elect) he believed that
three tests constituted a good yardstick by which to judge who might be saved: participation in the
two sacraments Baptism and the Lord’s Supper, an upright moral life, and a public profession of the
faith. These were adequate for a disciplined church on earth.

The consequence of faith to Calvin, far more than to Luther, is strenuous effort to introduce the
kingdom of God on earth. Though no longer judged by the law of God, the true Christian finds in the
law the divine pattern for moral character. Man is not justified by works, yet no justified man is
without works. No one can be a true Christian without aspiring to holiness in his life. This
rigorous pursuit of moral righteousness was one of the primary features of Calvinism. It made
character a fundamental test of genuine religious life and explains Calvinism’s dynamic social
activism. God calls the elect for his purpose!

Calvinism’s emphasis upon the sovereignty of God led in turn to a special view of the state. Luther
tended to consider the state supreme. The German princes often determined where and how the gospel
would be preached. But Calvin taught that no man, whether pope or king, has any claim to absolute
power. Calvin never preached the “right of revolution,” but he did encourage the growth of
representative assemblies and stressed their right to resist the tyranny of monarchs. Calvinist
resistance to the exercise of arbitrary power by monarchs was a key factor in the development of
modern constitutional governments.

The church, said Calvin, is not subject to secular government except in obviously secular matters.
On the other hand, the church has the obligation under the sovereign God, to guide the secular
authorities in spiritual matters. Such a vision sent Calvin’s followers throughout Europe as a
spiritual conspiracy seeking the overthrow of false religion and restrictive governments.

Many zealous disciples considered Geneva a beachhead established by God. It was a promise that the
kingdom would someday be organized. When they left Geneva they returned to their own countries to
establish Calvinist principles there. As a result, Calvinism rapidly assumed international
dimensions.

In France Calvinism remained a minority, but, thanks to influential converts among the nobility, the
movement gained an importance out of all proportion to its numbers. Known as Huguenots, French
Calvinists were threatening to seize leadership of the country when thousands of them were
ruthlessly massacred on St. Bartholomew’s Day in 1572. They remained a significant minority but
never again a serious challenge to the Catholic throne.

In the Netherlands Calvinism offered a rallying point for opposition to the oppressive rule of
Catholic Spain. Calvinist ministers were among the earliest leaders of resistance groups. Today we
would call them freedom fighters or perhaps guerrillas. The liberation leader of the national party
in the northern province of the Netherlands was William the Silent. He joined the reformed church in
1573 and during the next decade helped to fashion a Dutch republic. Today’s national anthem of the
Netherlands, “The Song of the Prince,” was written for William’s followers.

In Scotland the Calvinists created something unique in sixteenth-century Europe: a land of one
religion ruled by a monarch of another.

The monarch was Mary Queen of Scots, an eighteen-year-old girl living abroad. She married into the
French royal family, and the Scots as well as many Englishmen feared that she might deliver Scotland
to the French. One man, however, preached everywhere the notion that the people of Scotland could
challenge the rule of their queen. That man was John Knox.

\section{JOHN KNOX’S SCOTLAND}

Knox was a restless activist who had tried earlier to point England in the direction of Calvinism.
Like many others, however, he was forced to flee England overnight when, in 1553, the country
returned to the Catholic faith under Henry VIII’s daughter, Mary I. The queen’s persecution of
Protestant leaders earned her the title Bloody Mary.

Knox escaped to the Continent, where he developed the theory that Protestants had the right to
resist, by force if necessary, any Roman Catholic ruler who tried to prevent their worship and
mission. That was farther than Calvin himself was willing to go, but many of the nobles in Scotland
found the idea attractive.

When civil war broke out in Scotland in 1559, Knox rushed home. By the summer of 1560 the Calvinists
were in control of Edinburgh. Knox drafted the articles of religion that parliament accepted for the
country, thereby abolishing Roman Catholicism.

Next year when Mary Queen of Scots, now a nineteen-year-old widow, decided to return to her kingdom,
she found it in the lap of Protestant “heresy.” Over the next few years Knox, the passionate
preacher of Calvinism, and Mary, the young queen of Scotland, came to symbolize the Reformation
conflict: Protestant against Catholic, but also the democratic claims of Calvinism against the
monarchy’s power to appoint bishops. Events in Scotland moved in Knox’s direction. Even though
Mary’s descendants tried to turn back the clock, Scotland remained the most devoutly Calvinist
country in the world.

Thus, when Calvin died in 1564, he left far more than a reformed Geneva. All over Europe, and soon
in distant America, he had followers who were eager to continue the game he had entered that fateful
night in Farel’s rowdy and restless Geneva.

\readinghead{Suggestions for Further Reading}
\begin{readinglist}
\entry *Cottret, Bernard. Calvin: A Biography. Eerdmans, Grand Rapids, 2000.
\entry Ferm, Vergilius. Classics of Protestantism. New York: Philosophical Library, 1959.
\entry Harbison, E. Harris. The Christian Scholar in the Age of the Reformation. New York: Charles
  Scribner’s Sons, 1956.
\entry McGrath, Alister. A Life of John Calvin: A Study in the Shaping of Western Culture. Oxford:
  Basil Blackwell, 1990.
\entry McNeill, John T. The History and Character of Calvinism. New York: Oxford University Press,
  1967.
\entry Parker, T. H. L. Portrait of Calvin. Philadelphia: Westminster, John Knox Press, 2007.
\entry Walker, Williston. John Calvin, the Organizer of Reformed Protestantism. New York: Schocken,
  1969.
\end{readinglist}
